\documentclass[10pt,a4paper]{article}
\usepackage[T1]{fontenc}
\usepackage[utf8]{inputenc}
\usepackage{lmodern}
\usepackage{amsmath,amssymb}
\usepackage[margin=21mm]{geometry}
\usepackage{microtype}
\usepackage[hidelinks]{hyperref}
\usepackage{enumitem}
\setlist[itemize]{leftmargin=*,nosep,topsep=3pt}
\setlength{\parindent}{0pt}
\setlength{\parskip}{5pt}
\setlength{\emergencystretch}{3em}
\newcommand{\question}[1]{\par\medskip{\normalsize\bfseries #1}\par\smallskip}
\pagestyle{plain}

\begin{document}
{\Large\bfseries Neural signatures of hyperdirect pathway activity in Parkinson's disease\par}

\textbf{Paper:} Oswal et al. (2021).\quad Human recordings, tractography and modelling.\par
\href{https://doi.org/10.1038/s41467-021-25366-0}{DOI: \nolinkurl{10.1038/s41467-021-25366-0}}\par
{\footnotesize Page references below use PDF page numbers, including front matter.}

\question{To which general question does this paper relate?}
Network measurements and possible origins of beta synchronization. Helps distinguish frequency bands and pathways rather than treating beta as one homogeneous signal.

\question{How does it answer the general questions?}
\textbf{Measured.} STN/GPi recordings, MEG and tractography linked supplementary-motor-area--STN high-beta coherence (21--30~Hz) to hyperdirect fibre density. Directional analyses supported cortical drive at high beta; the corresponding relationship was absent for low beta (PDF pp.~1--7).

\textbf{Mechanism.} A coupled cortical/basal-ganglia model showed that enhanced high-beta hyperdirect input can provoke lower-beta subcortical synchronization. This is a model-based explanation supported by multimodal associations, not a selective clinical test proving that either band must be suppressed (PDF pp.~8--11).

\textbf{Model relevance.} Motivates adding cortical populations and hyperdirect coupling to an STN--GPe model. Transmission delays and time constants shape the simulated spectra. The full parameter set is referred to Supplementary Table~2; tract density and coherence should not be substituted directly for dynamical gains. The local main-article PDF alone does not supply a complete numerical notebook parameterization.

\question{What papers does it cite that could provide more information?}
\begin{itemize}
\item Nevado Holgado, Terry and Bogacz (2010), \emph{Conditions for the generation of beta oscillations in the subthalamic nucleus--globus pallidus network}: delayed rate model (ref.~33, PDF p.~13).
\item Hammond, Bergman and Brown (2007), \emph{Pathological synchronization in Parkinson's disease: networks, models and treatments} (ref.~38, PDF p.~13).
\end{itemize}

\end{document}
